\documentclass[11pt,a4paper]{article}

\usepackage[utf8]{inputenc}
\usepackage[T1]{fontenc}
\usepackage[french]{babel}
\usepackage[a4paper,margin=2cm]{geometry}
\usepackage{amsmath,amssymb,amsthm,mathtools}
\usepackage{lmodern}
\usepackage{enumitem}
\usepackage{xcolor}
\usepackage{hyperref}
\usepackage{fancyhdr}
\usepackage{tikz}
\usepackage{tcolorbox}
\usepackage{booktabs}
\usepackage{array}
\usepackage{titlesec}
\usepackage{physics}
\usepackage{multicol}
\usepackage{siunitx}

\hypersetup{
    colorlinks=true,
    linkcolor=black,
    urlcolor=blue
}

\setlength{\parindent}{0pt}
\setlength{\parskip}{0.5em}

\pagestyle{fancy}
\fancyhf{}
\fancyhead[L]{Spécialité Physique-Chimie -- Première générale}
\fancyhead[R]{Fiche de synthèse complète}
\fancyfoot[C]{\thepage}

\titleformat{\section}{\Large\bfseries}{\thesection.}{0.6em}{}
\titleformat{\subsection}{\large\bfseries}{\thesubsection.}{0.5em}{}
\titleformat{\subsubsection}{\normalsize\bfseries}{\thesubsubsection.}{0.5em}{}

\newtcolorbox{definitionbox}{
colback=blue!5!white,
colframe=blue!60!black,
title=\textbf{Définition},
fonttitle=\bfseries
}

\newtcolorbox{proprbox}{
colback=green!5!white,
colframe=green!50!black,
title=\textbf{Propriété / Résultat},
fonttitle=\bfseries
}

\newtcolorbox{sensbox}{
colback=orange!6!white,
colframe=orange!70!black,
title=\textbf{Sens physique},
fonttitle=\bfseries
}

\newtcolorbox{attentionbox}{
colback=red!4!white,
colframe=red!65!black,
title=\textbf{Attention / Erreur classique},
fonttitle=\bfseries
}

\newtcolorbox{terminalebox}{
colback=purple!5!white,
colframe=purple!60!black,
title=\textbf{Ouverture vers la Terminale},
fonttitle=\bfseries
}

\newtcolorbox{methodebox}{
colback=gray!8!white,
colframe=black!70,
title=\textbf{Méthode},
fonttitle=\bfseries
}

\newcommand{\R}{\mathbb{R}}
\newcommand{\C}{\mathbb{C}}

\begin{document}

\begin{center}
    {\LARGE \textbf{Fiche de synthèse complète}}\\[0.3em]
    {\Large \textbf{Physique-Chimie -- Spécialité Première générale}}\\[0.5em]
    {\large Notions, sens physique, liens entre chapitres, préparation à la Terminale}
\end{center}

\vspace{0.7em}

\tableofcontents

\newpage

\section{Comment lire un cours de physique-chimie}

\subsection{Pourquoi ce cours ne doit pas être appris comme une simple liste de formules}

En physique-chimie, une formule n'a de sens que si l'on comprend :
\begin{itemize}
    \item ce que représentent les grandeurs ;
    \item dans quelle situation on peut utiliser le modèle ;
    \item ce qu'on néglige ;
    \item ce que le résultat signifie physiquement.
\end{itemize}

Le but n'est donc pas seulement de \emph{savoir calculer}, mais surtout de \emph{comprendre le phénomène}.

\begin{sensbox}
La physique-chimie relie toujours :
\[
\text{observation} \longrightarrow \text{mesure} \longrightarrow \text{modèle} \longrightarrow \text{calcul} \longrightarrow \text{interprétation}
\]
\end{sensbox}

\subsection{Qu'est-ce qu'une définition ?}

\begin{definitionbox}
Une \textbf{définition} dit précisément ce qu'est une notion.  
Elle fixe le sens exact d'un mot scientifique.

Une définition ne se discute pas : elle sert de base au raisonnement.
\end{definitionbox}

\textbf{Exemples :}
\begin{itemize}
    \item La \textbf{concentration en quantité de matière} d'un soluté est le quotient de la quantité de matière de soluté par le volume de solution.
    \[
    c=\frac{n}{V}
    \]
    \item Une \textbf{onde mécanique progressive} est la propagation d'une perturbation dans un milieu matériel.
    \item Un \textbf{oxydant} est une espèce chimique capable de capter un ou plusieurs électrons.
\end{itemize}

\subsection{Qu'est-ce qu'une propriété ?}

\begin{proprbox}
Une \textbf{propriété} est un résultat vrai dans un certain cadre.  
Elle découle soit d'une définition, soit d'une expérience, soit d'un modèle admis.
\end{proprbox}

\textbf{Exemples :}
\begin{itemize}
    \item Si une onde périodique a pour période \(T\) et pour célérité \(v\), alors sa longueur d'onde vaut
    \[
    \lambda=vT.
    \]
    \item Si une transformation est totale, alors l'état final est déterminé par le réactif limitant.
    \item Dans le cas d'une force constante, le travail s'écrit
    \[
    W_{A\to B}(\vec F)=\vec F\cdot \overrightarrow{AB}.
    \]
\end{itemize}

\subsection{Qu'est-ce qu'une loi ?}

\begin{definitionbox}
Une \textbf{loi} est une relation générale utilisée pour modéliser un phénomène.
\end{definitionbox}

\textbf{Exemples :}
\begin{itemize}
    \item loi de Coulomb ;
    \item loi de Mariotte ;
    \item loi de Beer-Lambert ;
    \item relation entre débit de charges et intensité.
\end{itemize}

\subsection{Qu'est-ce qu'un modèle ?}

\begin{definitionbox}
Un \textbf{modèle} est une représentation simplifiée de la réalité qui permet d'expliquer, de prévoir ou de calculer.
\end{definitionbox}

\begin{sensbox}
Un modèle n'est pas la réalité entière.  
Il est choisi parce qu'il est assez simple pour travailler, mais assez fidèle pour répondre à la question posée.
\end{sensbox}

\textbf{Exemples :}
\begin{itemize}
    \item modéliser un objet par un point matériel ;
    \item modéliser une source réelle de tension par une source idéale et une résistance en série ;
    \item modéliser une transformation chimique par une équation de réaction ;
    \item modéliser la lumière tantôt comme onde, tantôt comme ensemble de photons.
\end{itemize}

\subsection{Grande idée du programme}

Tout le programme est construit autour de quelques idées transversales :
\begin{itemize}
    \item décrire un système ;
    \item suivre son évolution ;
    \item relier le microscopique et le macroscopique ;
    \item faire des bilans ;
    \item comprendre qu'un changement s'explique par une interaction, un transfert ou une conversion.
\end{itemize}

\newpage

\section{Constitution et transformations de la matière}

\subsection{Quantité de matière, masse molaire, concentration}

\subsubsection{La quantité de matière}

\begin{definitionbox}
La \textbf{quantité de matière} \(n\) mesure le nombre d'entités chimiques à l'échelle macroscopique.  
Elle s'exprime en mole (\(\si{mol}\)).
\end{definitionbox}

\begin{proprbox}
Pour un échantillon pur :
\[
n=\frac{m}{M}
\]
où \(m\) est la masse de l'échantillon et \(M\) la masse molaire.
\end{proprbox}

\begin{sensbox}
La mole est un outil de comptage.  
Comme les atomes et molécules sont beaucoup trop petits pour être comptés un par un, on les compte par ``paquets'' de très grande taille : les moles.
\end{sensbox}

\begin{attentionbox}
Ne pas confondre :
\begin{itemize}
    \item \textbf{masse} : grandeur mesurée en grammes ou kilogrammes ;
    \item \textbf{quantité de matière} : grandeur mesurée en moles.
\end{itemize}
Deux échantillons de même masse peuvent correspondre à des quantités de matière différentes.
\end{attentionbox}

\subsubsection{Masse molaire}

\begin{definitionbox}
La \textbf{masse molaire} \(M\) d'une espèce chimique est la masse d'une mole de cette espèce.
\end{definitionbox}

Pour une molécule, on l'obtient en additionnant les masses molaires atomiques des atomes qui la composent.

\textbf{Exemple :}
\[
M(\mathrm{H_2O})=2M(\mathrm H)+M(\mathrm O)
\]

\subsubsection{Concentration en quantité de matière}

\begin{definitionbox}
La \textbf{concentration molaire} d'un soluté en solution est
\[
c=\frac{n}{V}
\]
où \(n\) est la quantité de matière de soluté dissous et \(V\) le volume de solution.
\end{definitionbox}

\begin{sensbox}
La concentration mesure ``combien de soluté il y a par unité de volume''.  
Plus \(c\) est grande, plus la solution est riche en espèce dissoute.
\end{sensbox}

\begin{attentionbox}
Une solution n'est pas dite concentrée parce qu'elle contient ``beaucoup'' de matière en valeur absolue, mais parce qu'elle en contient beaucoup \emph{par litre}.
\end{attentionbox}

\subsubsection{Lien avec la Terminale}

\begin{terminalebox}
En Terminale, ces notions serviront partout :
\begin{itemize}
    \item en cinétique chimique ;
    \item en titrages plus avancés ;
    \item en chimie acide-base ;
    \item en électrochimie.
\end{itemize}
Si la notion de quantité de matière n'est pas solide, toute la chimie quantitative devient fragile.
\end{terminalebox}

\subsection{Absorbance, couleur, loi de Beer-Lambert}

\begin{definitionbox}
L'\textbf{absorbance} d'une solution mesure sa capacité à absorber une partie de la lumière qui la traverse.
\end{definitionbox}

\begin{sensbox}
Une solution colorée n'absorbe pas toutes les couleurs de la même façon.  
La couleur perçue correspond aux radiations transmises ou diffusées, et non à celles absorbées.
\end{sensbox}

\begin{proprbox}
Dans le domaine de validité du modèle de Beer-Lambert :
\[
A=k\,c
\]
où \(A\) est l'absorbance et \(c\) la concentration.
\end{proprbox}

\begin{methodebox}
Pour déterminer une concentration par étalonnage :
\begin{enumerate}
    \item on prépare une gamme étalon ;
    \item on mesure l'absorbance des solutions de référence ;
    \item on trace \(A=f(c)\) ;
    \item on mesure l'absorbance de la solution inconnue ;
    \item on lit ou calcule sa concentration.
\end{enumerate}
\end{methodebox}

\begin{terminalebox}
En Terminale, l'idée de relier un signal mesuré à une grandeur physique réapparaît dans beaucoup d'autres contextes : suivi cinétique, électricité, optique, radioactivité.
\end{terminalebox}

\subsection{Oxydants, réducteurs et réactions d'oxydo-réduction}

\begin{definitionbox}
Un \textbf{oxydant} est une espèce capable de capter un ou plusieurs électrons.  
Un \textbf{réducteur} est une espèce capable de céder un ou plusieurs électrons.
\end{definitionbox}

\begin{definitionbox}
Un \textbf{couple oxydant/réducteur} est l'association de deux formes chimiques reliées par gain ou perte d'électrons.
\end{definitionbox}

\begin{sensbox}
Une réaction d'oxydo-réduction est un échange d'électrons.  
C'est l'équivalent, pour les électrons, de ce que la réaction acide-base sera pour les protons.
\end{sensbox}

\textbf{Exemple de demi-équation :}
\[
\mathrm{Cu^{2+}+2e^- \rightarrow Cu}
\]

\begin{attentionbox}
Erreur classique : dire ``l'oxydant s'oxyde''.  
C'est faux. Un oxydant \textbf{oxyde l'autre espèce} et \textbf{se réduit lui-même}.
\end{attentionbox}

\subsection{Avancement et tableau d'avancement}

\begin{definitionbox}
L'\textbf{avancement} \(x\) d'une transformation chimique mesure son état de progression.
\end{definitionbox}

Pour une réaction
\[
aA+bB \rightarrow cC+dD,
\]
les quantités de matière évoluent selon :
\[
n_A=n_{A,0}-ax,\qquad
n_B=n_{B,0}-bx,\qquad
n_C=n_{C,0}+cx,\qquad
n_D=n_{D,0}+dx.
\]

\begin{sensbox}
L'avancement est un compteur de réaction.  
Il ne représente pas une espèce chimique en plus, mais la progression globale de la transformation.
\end{sensbox}

\begin{methodebox}
Pour construire un tableau d'avancement :
\begin{enumerate}
    \item écrire l'équation chimique ajustée ;
    \item déterminer les quantités initiales ;
    \item écrire l'évolution en fonction de \(x\) ;
    \item écrire l'état final ;
    \item utiliser le réactif limitant ou une donnée expérimentale pour trouver \(x_f\).
\end{enumerate}
\end{methodebox}

\begin{proprbox}
Dans le cas d'une transformation totale, l'état final est piloté par le réactif limitant.
\end{proprbox}

\begin{terminalebox}
En Terminale, l'idée d'avancement reviendra dans les bilans de matière, les suivis temporels, les titrages, et sera articulée avec des notions plus fines comme la cinétique et les évolutions non instantanées.
\end{terminalebox}

\subsection{Transformation totale, non totale, mélange stœchiométrique}

\begin{definitionbox}
Une transformation est dite \textbf{totale} si le réactif limitant est entièrement consommé.
\end{definitionbox}

\begin{definitionbox}
Un \textbf{mélange stœchiométrique} est un mélange dans lequel les réactifs sont introduits dans les proportions exactes imposées par les coefficients de l'équation chimique.
\end{definitionbox}

\begin{sensbox}
Un mélange stœchiométrique correspond à une situation idéale où aucun réactif n'est en excès si la transformation est totale.
\end{sensbox}

\subsection{Titrage colorimétrique}

\begin{definitionbox}
Un \textbf{titrage} permet de déterminer la quantité de matière ou la concentration d'une espèce chimique en la faisant réagir avec une espèce de concentration connue.
\end{definitionbox}

\begin{definitionbox}
L'\textbf{équivalence} est l'instant où les réactifs ont été introduits dans les proportions stœchiométriques.
\end{definitionbox}

\begin{proprbox}
À l'équivalence :
\[
\frac{n_A}{a}=\frac{n_B}{b}
\]
si \(a\) et \(b\) sont les coefficients stœchiométriques.
\end{proprbox}

\begin{sensbox}
Avant l'équivalence, un réactif titré est encore présent en excès.  
Après l'équivalence, c'est le titrant qui devient en excès.  
L'équivalence marque donc un \textbf{changement de réactif limitant}.
\end{sensbox}

\begin{attentionbox}
Il ne faut jamais utiliser la relation à l'équivalence sans avoir écrit d'abord la réaction support du titrage.
\end{attentionbox}

\subsection{Structure des entités, schémas de Lewis et géométrie}

\begin{definitionbox}
Le \textbf{schéma de Lewis} représente les doublets liants et non liants autour des atomes.
\end{definitionbox}

\begin{sensbox}
Le schéma de Lewis permet de relier la structure électronique à la forme de la molécule.  
On passe ainsi du microscopique (électrons, liaisons) au visible ou au mesurable (géométrie, polarité, solubilité).
\end{sensbox}

\textbf{Exemples de géométries :}
\begin{itemize}
    \item \(\mathrm{CO_2}\) : linéaire ;
    \item \(\mathrm{H_2O}\) : coudée ;
    \item \(\mathrm{CH_4}\) : tétraédrique.
\end{itemize}

\subsection{Électronégativité, polarité des liaisons et des molécules}

\begin{definitionbox}
L'\textbf{électronégativité} d'un atome mesure sa capacité à attirer vers lui les électrons d'une liaison.
\end{definitionbox}

\begin{definitionbox}
Une liaison covalente est \textbf{polarisée} si les deux atomes n'attirent pas les électrons de la même manière.
\end{definitionbox}

\begin{definitionbox}
Une molécule est \textbf{polaire} si la répartition des charges n'y est pas symétrique.
\end{definitionbox}

\begin{sensbox}
La polarité relie directement :
\[
\text{structure électronique} \longrightarrow \text{géométrie} \longrightarrow \text{propriétés macroscopiques}
\]
Elle explique la solubilité, les interactions entre molécules et donc beaucoup de propriétés physiques.
\end{sensbox}

\begin{attentionbox}
Une molécule peut posséder des liaisons polaires mais être globalement apolaire si la géométrie compense les polarités.
\end{attentionbox}

\subsection{Interactions entre entités, cohésion, solubilité, miscibilité}

\begin{definitionbox}
La \textbf{cohésion} d'une substance s'explique par les interactions entre ses entités microscopiques.
\end{definitionbox}

Types d'interactions à connaître :
\begin{itemize}
    \item interaction ion-ion ;
    \item interaction entre entités polaires ;
    \item interaction entre entités apolaires ;
    \item pont hydrogène.
\end{itemize}

\begin{sensbox}
La matière tient ensemble parce que ses entités interagissent.  
Comprendre les interactions, c'est comprendre pourquoi une substance est solide, liquide, soluble, peu soluble, miscible ou non miscible.
\end{sensbox}

\subsubsection{Dissolution d'un solide ionique}

\begin{definitionbox}
La dissolution d'un solide ionique dans l'eau correspond à la dissociation du cristal en ions hydratés.
\end{definitionbox}

\textbf{Exemple :}
\[
\mathrm{NaCl_{(s)} \rightarrow Na^+_{(aq)} + Cl^-_{(aq)}}
\]

\begin{sensbox}
L'eau stabilise les ions grâce à sa polarité.  
Elle arrache les ions au cristal et les entoure : c'est la \textbf{solvatation}.
\end{sensbox}

\subsubsection{Extraction par solvant}

\begin{definitionbox}
L'\textbf{extraction liquide-liquide} consiste à transférer une espèce chimique d'un solvant vers un autre dans lequel elle est plus soluble.
\end{definitionbox}

\begin{sensbox}
``Le semblable dissout le semblable'' : une espèce polaire se dissout mieux dans un solvant polaire, une espèce peu polaire dans un solvant peu polaire.
\end{sensbox}

\subsubsection{Savons et amphiphilie}

\begin{definitionbox}
Une espèce \textbf{amphiphile} possède à la fois une partie hydrophile et une partie lipophile.
\end{definitionbox}

\begin{sensbox}
Le savon peut interagir à la fois avec l'eau et avec les graisses.  
C'est cette double affinité qui explique ses propriétés lavantes.
\end{sensbox}

\subsection{Premières notions de chimie organique}

\subsubsection{Groupes caractéristiques}

\begin{definitionbox}
Un \textbf{groupe caractéristique} est un assemblage d'atomes responsable des propriétés chimiques principales d'une molécule organique.
\end{definitionbox}

Familles à connaître :
\begin{itemize}
    \item alcool ;
    \item aldéhyde ;
    \item cétone ;
    \item acide carboxylique.
\end{itemize}

\begin{sensbox}
En chimie organique, la chaîne carbonée donne une partie de l'identité de la molécule, mais c'est souvent le groupe caractéristique qui commande son comportement chimique.
\end{sensbox}

\subsubsection{Spectroscopie infrarouge}

\begin{definitionbox}
La spectroscopie infrarouge permet d'identifier des liaisons chimiques à partir des absorptions du rayonnement infrarouge.
\end{definitionbox}

\begin{sensbox}
Une liaison peut vibrer comme un ressort.  
Certaines fréquences lumineuses excitent ces vibrations : on observe alors des bandes caractéristiques dans le spectre.
\end{sensbox}

\subsubsection{Synthèse organique et rendement}

\begin{definitionbox}
Le \textbf{rendement} d'une synthèse est le rapport entre la quantité de produit réellement obtenue et la quantité maximale théorique.
\end{definitionbox}

\[
r=\frac{n_{\text{obtenu}}}{n_{\text{théorique}}}\times 100
\]

\begin{sensbox}
Le rendement mesure l'efficacité d'une synthèse.  
Il permet de comparer ce qu'on espérait obtenir à ce qu'on a réellement obtenu.
\end{sensbox}

\begin{attentionbox}
Un rendement inférieur à \(100\%\) peut venir :
\begin{itemize}
    \item d'une transformation incomplète ;
    \item de pertes lors de l'isolement ou de la purification ;
    \item de réactions parasites.
\end{itemize}
\end{attentionbox}

\subsubsection{Combustions et énergie de liaison}

\begin{definitionbox}
Une combustion est une transformation chimique exothermique au cours de laquelle un combustible réagit avec un comburant, généralement le dioxygène.
\end{definitionbox}

\begin{sensbox}
Au niveau microscopique, une réaction chimique correspond à des ruptures et formations de liaisons.  
L'énergie libérée ou absorbée vient du bilan entre ces ruptures et ces formations.
\end{sensbox}

\begin{terminalebox}
Cette idée prépare la Terminale en installant une manière énergétique de lire les transformations : on ne regarde plus seulement \emph{ce qui réagit}, mais aussi \emph{ce que cela coûte ou rapporte énergétiquement}.
\end{terminalebox}

\newpage

\section{Mouvement et interactions}

\subsection{Interactions fondamentales et champ}

\subsubsection{Charge électrique}

\begin{definitionbox}
La \textbf{charge électrique} est une propriété de certaines particules qui rend compte des interactions électrostatiques.
\end{definitionbox}

\begin{sensbox}
La charge joue, en électrostatique, un rôle analogue à celui de la masse en gravitation : c'est ce qui rend possible l'interaction.
\end{sensbox}

\subsubsection{Loi de Coulomb}

\begin{proprbox}
Deux charges ponctuelles exercent l'une sur l'autre des forces électrostatiques dont l'intensité dépend de leurs charges et de leur distance.
\end{proprbox}

\begin{sensbox}
Plus les charges sont grandes, plus l'interaction est forte.  
Plus elles sont éloignées, plus l'interaction s'affaiblit.
\end{sensbox}

\subsubsection{Gravitation}

\begin{proprbox}
Deux masses s'attirent mutuellement : c'est l'interaction gravitationnelle.
\end{proprbox}

\begin{sensbox}
L'intérêt pédagogique est de comparer gravitation et électrostatique : ce sont deux interactions à distance, décrites par des lois de forme voisine.
\end{sensbox}

\subsubsection{Champ}

\begin{definitionbox}
Un \textbf{champ} décrit en chaque point de l'espace l'influence qu'exercerait une source sur un objet-test placé en ce point.
\end{definitionbox}

\begin{sensbox}
Le champ permet de passer d'une vision ``objet contre objet'' à une vision plus locale : l'espace lui-même est décrit comme porteur d'une influence.
\end{sensbox}

\begin{terminalebox}
La notion de champ est fondamentale pour la Terminale.  
Elle réapparaît en gravitation, en électrostatique, dans les mouvements orbitaux, et plus largement dans toute la physique moderne.
\end{terminalebox}

\subsection{Fluide au repos}

\subsubsection{Grandeurs macroscopiques}

\begin{definitionbox}
Un fluide au repos est décrit macroscopiquement par des grandeurs comme :
\begin{itemize}
    \item la masse volumique ;
    \item la pression ;
    \item la température.
\end{itemize}
\end{definitionbox}

\begin{sensbox}
Le niveau microscopique parle de molécules en agitation ; le niveau macroscopique parle de grandeurs mesurables comme la pression ou la température.
\end{sensbox}

\subsubsection{Loi de Mariotte}

\begin{proprbox}
À température constante pour une quantité fixée de gaz :
\[
P\times V=\text{constante}
\]
\end{proprbox}

\begin{sensbox}
Si l'on diminue le volume disponible pour les particules d'un gaz, elles frappent plus souvent les parois : la pression augmente.
\end{sensbox}

\subsubsection{Forces pressantes}

\begin{proprbox}
Pour une surface plane \(S\) soumise à une pression \(P\), la valeur de la force pressante est :
\[
F=P\times S
\]
\end{proprbox}

\begin{sensbox}
La pression mesure une action répartie sur une surface.  
La force pressante traduit donc une action globale obtenue en combinant pression et surface.
\end{sensbox}

\subsubsection{Statique des fluides}

\begin{proprbox}
Dans un fluide incompressible au repos :
\[
P_2-P_1=\rho g(z_1-z_2)
\]
\end{proprbox}

\begin{sensbox}
Plus on descend dans un liquide, plus la pression augmente car le fluide situé au-dessus ``pèse'' sur les couches inférieures.
\end{sensbox}

\subsection{Mouvement d'un système}

\subsubsection{Vecteur vitesse et variation du vecteur vitesse}

\begin{definitionbox}
Le \textbf{vecteur vitesse} décrit à la fois la rapidité, la direction et le sens du mouvement.
\end{definitionbox}

\begin{definitionbox}
La \textbf{variation du vecteur vitesse} traduit comment le mouvement change entre deux instants.
\end{definitionbox}

\begin{sensbox}
Changer de vitesse ne signifie pas seulement aller plus vite ou moins vite : changer de direction, c'est déjà changer de vitesse au sens vectoriel.
\end{sensbox}

\subsubsection{Lien entre forces et variation de vitesse}

\begin{proprbox}
La somme des forces appliquées à un système est liée à la variation de son vecteur vitesse.
\end{proprbox}

\begin{sensbox}
Les forces ne commandent pas directement la vitesse, mais sa variation.  
C'est une idée centrale de toute la mécanique newtonienne.
\end{sensbox}

\begin{attentionbox}
Erreur fréquente : penser qu'une force est nécessaire pour maintenir un mouvement uniforme.  
En réalité, dans le cadre du principe d'inertie, une vitesse constante correspond à une résultante des forces nulle.
\end{attentionbox}

\begin{terminalebox}
En Terminale, ce lien sera reformulé de façon plus aboutie avec la deuxième loi de Newton dans sa forme différentielle et dans des mouvements plus riches, notamment circulaires et orbitaux.
\end{terminalebox}

\newpage

\section{Énergie : conversions et transferts}

\subsection{Aspects énergétiques des phénomènes électriques}

\subsubsection{Courant électrique et débit de charges}

\begin{definitionbox}
L'intensité du courant électrique mesure le débit de charges traversant une section de circuit.
\end{definitionbox}

\begin{proprbox}
\[
I=\frac{\Delta Q}{\Delta t}
\]
\end{proprbox}

\begin{sensbox}
Le courant n'est pas une réserve stockée : c'est un flux.  
Comme pour un débit d'eau, il dit combien de charge passe par seconde.
\end{sensbox}

\begin{attentionbox}
Ne pas confondre :
\begin{itemize}
    \item \textbf{charge} : quantité électrique ;
    \item \textbf{courant} : débit de charge.
\end{itemize}
\end{attentionbox}

\subsubsection{Source idéale et source réelle}

\begin{definitionbox}
Une source réelle de tension continue se modélise par une source idéale en série avec une résistance.
\end{definitionbox}

\begin{sensbox}
Aucun générateur réel n'est parfait.  
Le modèle de la source réelle rappelle qu'une partie de l'énergie fournie peut être perdue dans le générateur lui-même.
\end{sensbox}

\subsubsection{Puissance électrique, énergie, rendement}

\begin{definitionbox}
La \textbf{puissance} est la vitesse de transfert ou de conversion de l'énergie.
\end{definitionbox}

\begin{definitionbox}
Le \textbf{rendement} d'un convertisseur est le rapport entre la puissance utile et la puissance reçue.
\end{definitionbox}

\[
\eta=\frac{P_{\text{utile}}}{P_{\text{reçue}}}
\]

\begin{sensbox}
Un rendement mesure l'efficacité.  
Un rendement inférieur à \(1\) signifie qu'une partie de l'énergie reçue est dissipée, souvent par effet Joule.
\end{sensbox}

\begin{attentionbox}
Énergie et puissance sont souvent confondues :
\begin{itemize}
    \item l'énergie est une quantité ;
    \item la puissance est un rythme.
\end{itemize}
\end{attentionbox}

\subsection{Aspects énergétiques des phénomènes mécaniques}

\subsubsection{Énergie cinétique}

\begin{definitionbox}
L'\textbf{énergie cinétique} d'un système modélisé par un point matériel est l'énergie liée à son mouvement :
\[
E_c=\frac12 mv^2
\]
\end{definitionbox}

\begin{sensbox}
Plus un système est massif et rapide, plus son mouvement peut produire d'effet : choc, déformation, échauffement, freinage.
\end{sensbox}

\subsubsection{Travail d'une force}

\begin{definitionbox}
Le \textbf{travail} d'une force mesure le transfert d'énergie associé à cette force lors d'un déplacement.
\end{definitionbox}

\begin{proprbox}
Pour une force constante :
\[
W_{A\to B}(\vec F)=\vec F\cdot \overrightarrow{AB}
\]
\end{proprbox}

\begin{sensbox}
Le travail est la part de l'action de la force qui accompagne effectivement le déplacement.
\end{sensbox}

\subsubsection{Théorème de l'énergie cinétique}

\begin{proprbox}
La variation d'énergie cinétique d'un système est égale à la somme des travaux des forces appliquées :
\[
\Delta E_c=\sum W(\vec F)
\]
\end{proprbox}

\begin{sensbox}
Ce théorème donne une autre manière de relier forces et mouvement : non plus par l'accélération, mais par un bilan énergétique.
\end{sensbox}

\begin{terminalebox}
En Terminale, cette lecture énergétique devient encore plus puissante car elle permet d'étudier des systèmes complexes sans toujours revenir aux équations vectorielles du mouvement.
\end{terminalebox}

\subsubsection{Énergie potentielle de pesanteur}

\begin{definitionbox}
L'énergie potentielle de pesanteur d'un système près de la surface de la Terre s'écrit :
\[
E_p=mgz
\]
à une constante de référence près.
\end{definitionbox}

\begin{sensbox}
Elle mesure une réserve d'énergie liée à la position verticale.  
Mettre un objet plus haut, c'est lui donner la possibilité de redescendre en transformant cette énergie en autre chose, souvent en énergie cinétique.
\end{sensbox}

\subsubsection{Énergie mécanique}

\begin{definitionbox}
L'énergie mécanique est la somme de l'énergie cinétique et de l'énergie potentielle :
\[
E_m=E_c+E_p
\]
\end{definitionbox}

\begin{proprbox}
Si les forces non conservatives sont négligeables, l'énergie mécanique se conserve.
\end{proprbox}

\begin{sensbox}
L'énergie mécanique est le budget lié au mouvement et à la position.  
Quand il n'y a pas de dissipation, ce budget se redistribue mais ne change pas.
\end{sensbox}

\begin{attentionbox}
La conservation de l'énergie mécanique ne vaut pas toujours.  
Les frottements font généralement diminuer l'énergie mécanique : une partie est dissipée, souvent sous forme thermique.
\end{attentionbox}

\newpage

\section{Ondes et signaux}

\subsection{Ondes mécaniques}

\begin{definitionbox}
Une \textbf{onde mécanique progressive} est la propagation d'une perturbation dans un milieu matériel, sans transport global de matière.
\end{definitionbox}

\begin{sensbox}
Ce qui se propage, ce n'est pas la matière elle-même, mais une déformation, une vibration, une perturbation.
\end{sensbox}

\textbf{Exemples :}
\begin{itemize}
    \item onde sur une corde ;
    \item onde sonore ;
    \item houle ;
    \item ondes sismiques.
\end{itemize}

\begin{attentionbox}
Erreur classique : croire que la matière voyage avec l'onde.  
Sur une corde, la corde oscille localement ; elle ne part pas toute entière avec la perturbation.
\end{attentionbox}

\subsubsection{Célérité et retard}

\begin{definitionbox}
La \textbf{célérité} d'une onde est la vitesse de propagation de la perturbation.
\end{definitionbox}

\begin{proprbox}
\[
v=\frac{d}{\Delta t}
\]
\end{proprbox}

\begin{sensbox}
La célérité dépend du milieu.  
La même source peut produire une onde dont la célérité change si le milieu change.
\end{sensbox}

\subsubsection{Ondes périodiques, période, longueur d'onde}

\begin{definitionbox}
La \textbf{période} \(T\) est la durée d'un motif élémentaire.
\end{definitionbox}

\begin{definitionbox}
La \textbf{longueur d'onde} \(\lambda\) est la distance parcourue par l'onde pendant une période.
\end{definitionbox}

\begin{proprbox}
\[
\lambda=vT
\qquad \text{et} \qquad
f=\frac1T
\]
donc
\[
\lambda=\frac{v}{f}
\]
\end{proprbox}

\begin{sensbox}
Il faut distinguer :
\begin{itemize}
    \item la périodicité temporelle : ce qui se répète dans le temps en un point ;
    \item la périodicité spatiale : ce qui se répète dans l'espace à un instant donné.
\end{itemize}
\end{sensbox}

\begin{terminalebox}
Cette double lecture, temporelle et spatiale, prépare les chapitres plus riches de Terminale sur les ondes, les phénomènes de propagation et les interprétations fines des signaux.
\end{terminalebox}

\subsection{Optique géométrique : lentilles et images}

\subsubsection{Image par une lentille convergente}

\begin{definitionbox}
Une lentille mince convergente peut former l'image d'un objet en faisant converger ou diverger les rayons lumineux selon la position de l'objet.
\end{definitionbox}

\begin{sensbox}
L'optique géométrique est un modèle simple qui remplace la lumière par des rayons.  
Ce modèle est très utile pour comprendre la formation des images, même s'il ne décrit pas tout.
\end{sensbox}

\textbf{Notions à connaître :}
\begin{itemize}
    \item image réelle ;
    \item image virtuelle ;
    \item image droite ;
    \item image renversée ;
    \item distance focale ;
    \item grandissement.
\end{itemize}

\begin{proprbox}
On exploite les relations de conjugaison et de grandissement fournies par le cours.
\end{proprbox}

\begin{terminalebox}
En Terminale, l'optique sera enrichie par une lecture plus ondulatoire de la lumière.  
La Première installe donc un double regard : la lumière peut être modélisée géométriquement, mais aussi ondulatoirement et particulairement.
\end{terminalebox}

\subsection{Couleurs, synthèse additive et soustractive}

\begin{definitionbox}
La \textbf{synthèse additive} correspond à la superposition de lumières colorées.
\end{definitionbox}

\begin{definitionbox}
La \textbf{synthèse soustractive} correspond à la sélection de certaines couleurs par absorption, par exemple avec des filtres ou des pigments.
\end{definitionbox}

\begin{sensbox}
La couleur perçue dépend :
\begin{itemize}
    \item de la lumière reçue par l'objet ;
    \item de ce que l'objet absorbe ;
    \item de ce qu'il diffuse ou transmet.
\end{itemize}
La couleur n'est donc pas une propriété totalement indépendante de l'éclairage.
\end{sensbox}

\subsection{Lumière : modèles ondulatoire et particulaire}

\subsubsection{Onde électromagnétique}

\begin{definitionbox}
La lumière est une onde électromagnétique.
\end{definitionbox}

\begin{proprbox}
\[
\lambda=\frac{c}{\nu}
\]
où \(\lambda\) est la longueur d'onde, \(c\) la célérité de la lumière dans le vide, et \(\nu\) la fréquence.
\end{proprbox}

\begin{sensbox}
Cette relation montre qu'une même lumière peut être décrite par sa fréquence ou par sa longueur d'onde.
\end{sensbox}

\subsubsection{Photon}

\begin{definitionbox}
Le \textbf{photon} est le quantum d'énergie associé à la lumière.
\end{definitionbox}

\begin{proprbox}
\[
E=h\nu
\]
\end{proprbox}

\begin{sensbox}
Le modèle particulaire rappelle que la lumière échange son énergie de manière discrète avec la matière.
\end{sensbox}

\subsubsection{Spectres et niveaux d'énergie}

\begin{definitionbox}
Les atomes possèdent des niveaux d'énergie quantifiés.
\end{definitionbox}

\begin{sensbox}
Quand un atome absorbe ou émet de la lumière, il change de niveau d'énergie.  
La lumière échangée porte exactement la différence d'énergie entre ces niveaux.
\end{sensbox}

\begin{proprbox}
\[
\Delta E=h\nu
\qquad \text{et} \qquad
\lambda=\frac{c}{\nu}
\]
\end{proprbox}

\begin{terminalebox}
C'est une préparation directe à la Terminale : spectres, photons, interaction lumière-matière, quantification de l'énergie sont des idées centrales pour la suite.
\end{terminalebox}

\newpage

\section{Mesures, incertitudes et démarche scientifique}

\subsection{Pourquoi une mesure n'est jamais parfaite}

\begin{definitionbox}
Une \textbf{mesure} n'est jamais la valeur vraie absolue d'une grandeur : c'est une estimation.
\end{definitionbox}

\begin{sensbox}
Mesurer, ce n'est pas atteindre une exactitude parfaite ; c'est produire une valeur raisonnable assortie d'un degré de confiance.
\end{sensbox}

\subsection{Incertitude-type}

\begin{definitionbox}
L'\textbf{incertitude-type} traduit la dispersion ou l'imprécision associée à une mesure.
\end{definitionbox}

Deux approches :
\begin{itemize}
    \item \textbf{type A} : approche statistique à partir d'une série de mesures ;
    \item \textbf{type B} : approche à partir des caractéristiques de l'instrument ou du protocole.
\end{itemize}

\begin{attentionbox}
Donner une incertitude ne signifie pas que la mesure est mauvaise.  
Au contraire, c'est ce qui la rend scientifiquement exploitable.
\end{attentionbox}

\subsection{Chiffres significatifs}

\begin{definitionbox}
Le nombre de chiffres significatifs doit être cohérent avec la précision de la mesure.
\end{definitionbox}

\begin{sensbox}
Écrire trop de chiffres donne une illusion de précision.  
Un bon scientifique n'écrit pas plus de chiffres qu'il n'en connaît réellement.
\end{sensbox}

\subsection{Compétences de la démarche scientifique}

Les grandes compétences à travailler sont :
\begin{itemize}
    \item \textbf{s'approprier} ;
    \item \textbf{analyser / raisonner} ;
    \item \textbf{réaliser} ;
    \item \textbf{valider} ;
    \item \textbf{communiquer}.
\end{itemize}

\begin{sensbox}
Ces compétences ne sont pas séparées du cours : elles sont la manière concrète de faire de la physique-chimie.
\end{sensbox}

\newpage

\section{Grandes liaisons entre les notions}

\subsection{Le programme forme un tout}

Voici quelques liens fondamentaux :

\begin{itemize}
    \item \textbf{Structure microscopique \(\rightarrow\) propriétés macroscopiques} : polarité, solubilité, cohésion, couleur.
    \item \textbf{Forces \(\rightarrow\) mouvement} : interaction, champ, variation de vitesse.
    \item \textbf{Forces \(\rightarrow\) énergie} : travail, théorème de l'énergie cinétique.
    \item \textbf{Position \(\rightarrow\) énergie potentielle} : hauteur, champ de pesanteur, conservation.
    \item \textbf{Signal mesuré \(\rightarrow\) grandeur recherchée} : absorbance, titrage, intensité, spectre.
    \item \textbf{Périodicité temporelle \(\leftrightarrow\) périodicité spatiale} : onde, fréquence, longueur d'onde.
    \item \textbf{Modèle choisi \(\rightarrow\) résultat obtenu} : point matériel, source réelle, optique géométrique, photon.
\end{itemize}

\begin{sensbox}
Ce qu'il faut retenir, ce n'est pas une collection de chapitres isolés, mais une manière de penser :
\begin{center}
\textbf{décrire un système, choisir un modèle, relier les grandeurs, faire un bilan, interpréter.}
\end{center}
\end{sensbox}

\section{Ouvertures générales vers la Terminale}

\begin{terminalebox}
La classe de Première prépare la Terminale de plusieurs façons :

\begin{itemize}
    \item \textbf{en chimie} : bilans de matière, oxydoréduction, titrages, structure et énergie des transformations ;
    \item \textbf{en mécanique} : champ, forces, variation de vitesse, lecture énergétique ;
    \item \textbf{en électricité} : puissance, rendement, modèles de dipôles, transfert d'énergie ;
    \item \textbf{en ondes et en optique} : périodicité, longueur d'onde, spectres, photons, interaction lumière-matière ;
    \item \textbf{en méthode} : modéliser, mesurer, critiquer, comparer, raisonner avec cohérence.
\end{itemize}

Autrement dit, la Première installe les \textbf{idées fondatrices}, et la Terminale les \textbf{approfondit}, les \textbf{raffine} et les \textbf{généralise}.
\end{terminalebox}

\section{Résumé final : ce qu'il faut avoir compris}

\begin{itemize}
    \item Une \textbf{définition} dit ce qu'est une notion.
    \item Une \textbf{propriété} donne un résultat valable dans un cadre donné.
    \item Une \textbf{loi} modélise quantitativement un phénomène.
    \item Un \textbf{modèle} simplifie la réalité pour pouvoir raisonner.
    \item La physique-chimie ne consiste pas à réciter des formules, mais à comprendre :
    \[
    \text{quoi décrire, avec quelles grandeurs, dans quel modèle, et pour dire quoi du réel.}
    \]
\end{itemize}

\begin{center}
\Large \textbf{Fin de la fiche}
\end{center}

\end{document}